\section{Introduction}
\label{sec:introduction}

CXL-attached memory offers a practical way to expand server capacity without
populating every processor memory channel.  A host can access a Type-3 device
through ordinary loads and stores, and future fabrics can pool capacity across
machines~\cite{cxl30}.  The capacity is useful, but not free: additional hops,
controllers, and media increase access latency and constrain bandwidth.  The
operating system must therefore decide which pages remain in scarce local DRAM
(\fast) and which pages reside in larger CXL memory (\slow).

Existing tiering systems make this decision from recent page temperature.
They periodically sample accesses, classify hot pages, and migrate them toward
the fast tier~\cite{agarwal2017thermostat,yan2019nimble,raybuck2021hemem,
maruf2023tpp}.  This reactive loop is effective for long, stationary phases.
It is less effective for services whose memory demand changes at request or
pipeline boundaries.  By the time a burst has generated enough samples to
classify its pages as hot, the burst may already have completed on the slow
tier.  Average throughput can remain healthy even while the few requests that
encounter such transitions dominate tail latency~\cite{dean2013tail}.

This paper asks whether the kernel can act \emph{before} the first expensive
miss without requiring programmers to mark phases or objects.  Our key
observation is that many services revisit execution contexts---parser loops,
index probes, compaction workers, and model operators---before revisiting the
same groups of pages.  Sampled instruction pointers thus provide an early but
noisy phase signal.  A compact history can associate that signal with a
working-set fingerprint and pre-promote a small set of pages.  The mechanism
must nevertheless be conservative: an inaccurate predictor can evict the
foreground working set or saturate the migration path, producing precisely the
tail spikes it is intended to avoid.

We introduce \system, phase-aware memory tiering for latency-sensitive
services.  \system runs inside the Linux memory-management path and operates at
the granularity of a cgroup.  It observes sampled instruction pointers and
accessed pages, learns recurring phase signatures, and predicts the next
working set.  A migration-credit controller admits speculative promotions only
when the expected value exceeds their estimated interference cost.  Finally, a
deadline-aware reclaim policy lends local DRAM to predicted pages for a bounded
lease, preventing stale predictions from becoming permanent allocations.

The paper makes three contributions:

\begin{itemize}
  \item We identify \emph{phase exposure}---the interval between recognizable
        execution and the first slow-tier access---as a resource that reactive
        page-temperature policies leave unused.
  \item We design a phase predictor and credit-bounded promotion mechanism that
        converts this interval into early migrations while preserving a safe
        reactive fallback.
  \item We implement the design in Linux and evaluate latency, throughput,
        bandwidth, fairness, and prediction-error sensitivity on eight
        application mixes.  \system lowers \ptail latency by up to 58\% and bounds
        wasted migrations to 6.2\% of available copy bandwidth.
\end{itemize}

\paragraph{Scope.}
Our prototype manages volatile, cache-coherent expansion memory.  It does not
provide persistence semantics, device-level wear management, or fabric-wide
admission control.  The evaluation uses a calibrated NUMA emulation because
the tested CXL device does not expose programmable latency; \cref{sec:limits}
discusses the implications.
